% prezentace.tex -- Lekce 04: Cykly a animace
% preamble-prez.tex -- sdílená preamble pro prezentace (beamer)
% Includuje se z ucitel/lekceXX/prezentace.tex jako:
%   % preamble-prez.tex -- sdílená preamble pro prezentace (beamer)
% Includuje se z ucitel/lekceXX/prezentace.tex jako:
%   % preamble-prez.tex -- sdílená preamble pro prezentace (beamer)
% Includuje se z ucitel/lekceXX/prezentace.tex jako:
%   \input{../spolecne/preamble-prez.tex}

\documentclass[aspectratio=169]{beamer}

% --- Jazyk a font ---
\usepackage{polyglossia}
\setmainlanguage{czech}
\usepackage{fontspec}

% --- Téma ---
\usetheme{default}
\usecolortheme{default}

\definecolor{microbit-blue}{HTML}{1E4D8C}
\definecolor{microbit-lightblue}{HTML}{D6E4F7}
\definecolor{code-bg}{HTML}{F0F0F0}

\setbeamercolor{frametitle}{bg=microbit-blue, fg=white}
\setbeamercolor{title}{fg=microbit-blue}
\setbeamercolor{block title}{bg=microbit-blue, fg=white}
\setbeamercolor{block body}{bg=microbit-lightblue}
\setbeamercolor{itemize item}{fg=microbit-blue}
\setbeamercolor{itemize subitem}{fg=microbit-blue!70}
\setbeamerfont{frametitle}{size=\large, series=\bfseries}
\setbeamertemplate{navigation symbols}{}
\setbeamertemplate{footline}{%
  \leavevmode%
  \hbox{%
    \begin{beamercolorbox}[wd=.5\paperwidth,ht=2.5ex,dp=1.125ex,leftskip=6pt]{author in head/foot}%
      \usebeamerfont{author in head/foot}\textcolor{gray}{\small Úvod do Pythonu na Micro:bitu}
    \end{beamercolorbox}%
    \begin{beamercolorbox}[wd=.5\paperwidth,ht=2.5ex,dp=1.125ex,rightskip=6pt,right]{date in head/foot}%
      \usebeamerfont{date in head/foot}\textcolor{gray}{\small\insertframenumber\,/\,\inserttotalframenumber}
    \end{beamercolorbox}}%
}

% --- Česky korektní uvozovky ---
\usepackage{csquotes}
\newcommand{\uv}[1]{\enquote{#1}}

% --- Zdrojový kód (Python) ---
\usepackage{listings}
\lstset{
  language=Python,
  basicstyle=\ttfamily\small,
  backgroundcolor=\color{code-bg},
  frame=single,
  framerule=0pt,
  xleftmargin=6pt,
  xrightmargin=6pt,
  breaklines=true,
  showstringspaces=false,
  keywordstyle=\color{microbit-blue}\bfseries,
  stringstyle=\color{teal},
  commentstyle=\color{gray}\itshape,
  literate=
    {á}{{\'a}}1 {é}{{\'e}}1 {í}{{\'i}}1 {ó}{{\'o}}1 {ú}{{\'u}}1
    {ů}{{\r{u}}}1 {č}{{\v{c}}}1 {š}{{\v{s}}}1 {ž}{{\v{z}}}1
    {Á}{{\'A}}1 {É}{{\'E}}1 {Í}{{\'I}}1 {Ó}{{\'O}}1 {Ú}{{\'U}}1
    {Č}{{\v{C}}}1 {Š}{{\v{S}}}1 {Ž}{{\v{Z}}}1 {ň}{{\v{n}}}1 {ř}{{\v{r}}}1,
}

% --- Zvýraznění pojmů ---
\newcommand{\pojem}[1]{\textbf{\textcolor{microbit-blue}{#1}}}
\newcommand{\pyinline}[1]{\texttt{#1}}

% --- Grafika a diagramy ---
\usepackage{tikz}

% --- Ostatní ---
\usepackage{graphicx}
\usepackage{hyperref}
\hypersetup{colorlinks=true, urlcolor=microbit-blue}


\documentclass[aspectratio=169]{beamer}

% --- Jazyk a font ---
\usepackage{polyglossia}
\setmainlanguage{czech}
\usepackage{fontspec}

% --- Téma ---
\usetheme{default}
\usecolortheme{default}

\definecolor{microbit-blue}{HTML}{1E4D8C}
\definecolor{microbit-lightblue}{HTML}{D6E4F7}
\definecolor{code-bg}{HTML}{F0F0F0}

\setbeamercolor{frametitle}{bg=microbit-blue, fg=white}
\setbeamercolor{title}{fg=microbit-blue}
\setbeamercolor{block title}{bg=microbit-blue, fg=white}
\setbeamercolor{block body}{bg=microbit-lightblue}
\setbeamercolor{itemize item}{fg=microbit-blue}
\setbeamercolor{itemize subitem}{fg=microbit-blue!70}
\setbeamerfont{frametitle}{size=\large, series=\bfseries}
\setbeamertemplate{navigation symbols}{}
\setbeamertemplate{footline}{%
  \leavevmode%
  \hbox{%
    \begin{beamercolorbox}[wd=.5\paperwidth,ht=2.5ex,dp=1.125ex,leftskip=6pt]{author in head/foot}%
      \usebeamerfont{author in head/foot}\textcolor{gray}{\small Úvod do Pythonu na Micro:bitu}
    \end{beamercolorbox}%
    \begin{beamercolorbox}[wd=.5\paperwidth,ht=2.5ex,dp=1.125ex,rightskip=6pt,right]{date in head/foot}%
      \usebeamerfont{date in head/foot}\textcolor{gray}{\small\insertframenumber\,/\,\inserttotalframenumber}
    \end{beamercolorbox}}%
}

% --- Česky korektní uvozovky ---
\usepackage{csquotes}
\newcommand{\uv}[1]{\enquote{#1}}

% --- Zdrojový kód (Python) ---
\usepackage{listings}
\lstset{
  language=Python,
  basicstyle=\ttfamily\small,
  backgroundcolor=\color{code-bg},
  frame=single,
  framerule=0pt,
  xleftmargin=6pt,
  xrightmargin=6pt,
  breaklines=true,
  showstringspaces=false,
  keywordstyle=\color{microbit-blue}\bfseries,
  stringstyle=\color{teal},
  commentstyle=\color{gray}\itshape,
  literate=
    {á}{{\'a}}1 {é}{{\'e}}1 {í}{{\'i}}1 {ó}{{\'o}}1 {ú}{{\'u}}1
    {ů}{{\r{u}}}1 {č}{{\v{c}}}1 {š}{{\v{s}}}1 {ž}{{\v{z}}}1
    {Á}{{\'A}}1 {É}{{\'E}}1 {Í}{{\'I}}1 {Ó}{{\'O}}1 {Ú}{{\'U}}1
    {Č}{{\v{C}}}1 {Š}{{\v{S}}}1 {Ž}{{\v{Z}}}1 {ň}{{\v{n}}}1 {ř}{{\v{r}}}1,
}

% --- Zvýraznění pojmů ---
\newcommand{\pojem}[1]{\textbf{\textcolor{microbit-blue}{#1}}}
\newcommand{\pyinline}[1]{\texttt{#1}}

% --- Grafika a diagramy ---
\usepackage{tikz}

% --- Ostatní ---
\usepackage{graphicx}
\usepackage{hyperref}
\hypersetup{colorlinks=true, urlcolor=microbit-blue}


\documentclass[aspectratio=169]{beamer}

% --- Jazyk a font ---
\usepackage{polyglossia}
\setmainlanguage{czech}
\usepackage{fontspec}

% --- Téma ---
\usetheme{default}
\usecolortheme{default}

\definecolor{microbit-blue}{HTML}{1E4D8C}
\definecolor{microbit-lightblue}{HTML}{D6E4F7}
\definecolor{code-bg}{HTML}{F0F0F0}

\setbeamercolor{frametitle}{bg=microbit-blue, fg=white}
\setbeamercolor{title}{fg=microbit-blue}
\setbeamercolor{block title}{bg=microbit-blue, fg=white}
\setbeamercolor{block body}{bg=microbit-lightblue}
\setbeamercolor{itemize item}{fg=microbit-blue}
\setbeamercolor{itemize subitem}{fg=microbit-blue!70}
\setbeamerfont{frametitle}{size=\large, series=\bfseries}
\setbeamertemplate{navigation symbols}{}
\setbeamertemplate{footline}{%
  \leavevmode%
  \hbox{%
    \begin{beamercolorbox}[wd=.5\paperwidth,ht=2.5ex,dp=1.125ex,leftskip=6pt]{author in head/foot}%
      \usebeamerfont{author in head/foot}\textcolor{gray}{\small Úvod do Pythonu na Micro:bitu}
    \end{beamercolorbox}%
    \begin{beamercolorbox}[wd=.5\paperwidth,ht=2.5ex,dp=1.125ex,rightskip=6pt,right]{date in head/foot}%
      \usebeamerfont{date in head/foot}\textcolor{gray}{\small\insertframenumber\,/\,\inserttotalframenumber}
    \end{beamercolorbox}}%
}

% --- Česky korektní uvozovky ---
\usepackage{csquotes}
\newcommand{\uv}[1]{\enquote{#1}}

% --- Zdrojový kód (Python) ---
\usepackage{listings}
\lstset{
  language=Python,
  basicstyle=\ttfamily\small,
  backgroundcolor=\color{code-bg},
  frame=single,
  framerule=0pt,
  xleftmargin=6pt,
  xrightmargin=6pt,
  breaklines=true,
  showstringspaces=false,
  keywordstyle=\color{microbit-blue}\bfseries,
  stringstyle=\color{teal},
  commentstyle=\color{gray}\itshape,
  literate=
    {á}{{\'a}}1 {é}{{\'e}}1 {í}{{\'i}}1 {ó}{{\'o}}1 {ú}{{\'u}}1
    {ů}{{\r{u}}}1 {č}{{\v{c}}}1 {š}{{\v{s}}}1 {ž}{{\v{z}}}1
    {Á}{{\'A}}1 {É}{{\'E}}1 {Í}{{\'I}}1 {Ó}{{\'O}}1 {Ú}{{\'U}}1
    {Č}{{\v{C}}}1 {Š}{{\v{S}}}1 {Ž}{{\v{Z}}}1 {ň}{{\v{n}}}1 {ř}{{\v{r}}}1,
}

% --- Zvýraznění pojmů ---
\newcommand{\pojem}[1]{\textbf{\textcolor{microbit-blue}{#1}}}
\newcommand{\pyinline}[1]{\texttt{#1}}

% --- Grafika a diagramy ---
\usepackage{tikz}

% --- Ostatní ---
\usepackage{graphicx}
\usepackage{hyperref}
\hypersetup{colorlinks=true, urlcolor=microbit-blue}


\title{Lekce 04: Cykly a animace}
\subtitle{Úvod do Pythonu na Micro:bitu}
\date{}

\begin{document}

\begin{frame}
  \titlepage
\end{frame}

% ------------------------------------------------------------------
\begin{frame}[fragile]{\texttt{while True:} — teď pořádně}

  V minulé lekci jsme použili \pyinline{while True:} jako zaklínadlo.
  Teď víme, co to znamená:

  \bigskip
  \begin{block}{while cyklus}
    \begin{lstlisting}
while podminka:
    # opakuj, dokud podminka plati
    \end{lstlisting}
  \end{block}

  \bigskip
  \pyinline{True} je vždy pravda → cyklus \textbf{nikdy neskončí}.

  \medskip
  Používáme ho pro programy, které mají běžet pořád:\\
  sledovat tlačítka, senzory, zobrazovat animace.

  \bigskip
  \textcolor{gray}{\small Zastavení: tlačítko \textbf{Stop} v Thonny nebo \textbf{Reset} na micro:bitu.}
\end{frame}

% ------------------------------------------------------------------
\begin{frame}[fragile]{\texttt{sleep()} — pauza}

  \pyinline{sleep(ms)} pozastaví program na zadaný počet \textbf{milisekund}:

  \bigskip
  \begin{center}
  \begin{tabular}{ll}
    \toprule
    \textbf{Zápis} & \textbf{Pauza} \\
    \midrule
    \pyinline{sleep(100)}  & 0,1 sekundy \\
    \pyinline{sleep(500)}  & půl sekundy \\
    \pyinline{sleep(1000)} & 1 sekunda \\
    \pyinline{sleep(2000)} & 2 sekundy \\
    \bottomrule
  \end{tabular}
  \end{center}

  \bigskip
  \begin{lstlisting}
while True:
    display.show(Image.HEART)
    sleep(500)
    display.clear()
    sleep(500)
  \end{lstlisting}
\end{frame}

% ------------------------------------------------------------------
\begin{frame}[fragile]{\texttt{for} cyklus — opakuj N-krát}

  \pyinline{while True:} běží donekonečna.
  \pyinline{for} cyklus se zastaví po přesném počtu opakování:

  \bigskip
  \begin{lstlisting}
for i in range(5):
    display.show(str(i))
    sleep(600)
  \end{lstlisting}

  \bigskip
  \begin{itemize}
    \item \pyinline{range(5)} vygeneruje čísla \pyinline{0, 1, 2, 3, 4}
    \item Proměnná \pyinline{i} přijímá postupně každou hodnotu
    \item Po posledním opakování cyklus skončí a program pokračuje dál
  \end{itemize}

  \pause\bigskip
  \begin{alertblock}{Pozor: range začíná od nuly!}
    \pyinline{range(5)} → 0, 1, 2, 3, 4 \quad (ne 1 až 5)
  \end{alertblock}
\end{frame}

% ------------------------------------------------------------------
\begin{frame}[fragile]{Iterační proměnná \texttt{i}}

  Proměnnou \pyinline{i} lze použít přímo ve výpočtu nebo podmínce:

  \begin{lstlisting}
for i in range(5):
    display.set_pixel(i, 2, 9)   # sviti pixel na x=i, y=2
    sleep(200)
  \end{lstlisting}

  \bigskip
  \begin{center}
  \begin{tikzpicture}[scale=0.7]
    % 5x5 grid
    \foreach \x in {0,...,4}
      \foreach \y in {0,...,4}
        \draw[gray!30] (\x*0.8, -\y*0.8) rectangle (\x*0.8+0.7, -\y*0.8+0.7);
    % middle row highlight (y=2 -> row index 2 from top)
    \foreach \x in {0,...,4}
      \draw[microbit-blue, very thick] (\x*0.8, -2*0.8) rectangle (\x*0.8+0.7, -2*0.8+0.7);
    % animated pixel (show i=2 as example)
    \fill[microbit-blue] (2*0.8+0.1, -2*0.8+0.1) rectangle (2*0.8+0.6, -2*0.8+0.6);
    % labels
    \foreach \x in {0,...,4}
      \node[font=\tiny, gray] at (\x*0.8+0.35, 0.5) {\x};
    \node[font=\small, microbit-blue] at (-0.6, -2*0.8+0.35) {y=2};
  \end{tikzpicture}
  \end{center}

  \textcolor{gray}{\small \pyinline{i} postupně nabývá hodnot 0, 1, 2, 3, 4 → pixel se pohybuje.}
\end{frame}

% ------------------------------------------------------------------
\begin{frame}[fragile]{Animace — sekvence obrázků}

  Obrázky uložíme do \textbf{seznamu} a procházíme je \pyinline{for} cyklem:

  \begin{lstlisting}
from microbit import *

snimky = [Image.HAPPY, Image.SMILE, Image.SAD]

while True:
    for obrazek in snimky:
        display.show(obrazek)
        sleep(400)
  \end{lstlisting}

  \bigskip
  \begin{itemize}
    \item Proměnná \pyinline{obrazek} přijímá postupně každý prvek seznamu
    \item \pyinline{while True:} zajistí, že animace běží stále dokola
    \item Délku pauzy (400 ms) lze libovolně měnit
  \end{itemize}

  \pause\bigskip
  \begin{exampleblock}{Vestavěná animace}
    \pyinline{display.show(Image.ALL\_CLOCKS, delay=100, loop=True)}
  \end{exampleblock}
\end{frame}

% ------------------------------------------------------------------
\begin{frame}{Co jsme se naučili}
  \begin{itemize}
    \item[\checkmark] \pyinline{while True:} — nekonečný cyklus; typický pro senzory a animace
    \item[\checkmark] \pyinline{sleep(ms)} — pauza v milisekundách
    \item[\checkmark] \pyinline{display.clear()} — zhasne celý displej
    \item[\checkmark] \pyinline{for i in range(n):} — opakuj přesně n-krát
    \item[\checkmark] \pyinline{range(n)} začíná od 0, končí n-1
    \item[\checkmark] Iterační proměnná \pyinline{i} lze použít ve výpočtu
    \item[\checkmark] Animace = seznam obrázků procházený \pyinline{for} cyklem
    \item[\checkmark] \pyinline{display.set_pixel(x, y, jas)} — ovládání jedné LED
  \end{itemize}
\end{frame}

\end{document}
